\section{Reading attention where internals are needed}
\label{sec:attention}

\begin{figure}[t]
\centering
\includegraphics[width=\columnwidth]{fig3_attention}
\caption{The attention tier on every powered run: the length floor, the head-averaged spectral profile, the per-head profile, LapEigvals and the token-role probe. Keeping heads apart moves the attention readout from the floor to within a few points of the probe.}
\label{fig:attention}
\end{figure}

An operator whose stack exposes attention weights but not activations wants to know how much of the probe an attention-only readout recovers. On the runs where internals are needed, the per-head spectral profile and LapEigvals trail the token-role probe by a few points (\cref{fig:attention}): \GapProbePerHeadLlamaOneBBfcl{} on Llama-3.2-1B and \GapProbePerHeadGemmaBfcl{} on Gemma-3-1B on BFCL, with the two attention readouts within noise of each other on every run (per-head minus LapEigvals \GapPerHeadLapEigLlamaOneBBfcl{} and \GapPerHeadLapEigGemmaBfcl{}). SinkProbe and the anchored readout of the call's own rows sit in the same band. On the runs where confidence suffices every attention readout sits below confidence, as the probe does.

What decides whether an attention readout works at all is whether the heads are averaged before anything is computed. The head-averaged spectral profile, which is how most of this literature reads a layer, does not clear the length floor on several runs, and keeping heads apart is worth \GapPerHeadAvgLlamaOneBBfcl{} on Llama-3.2-1B on BFCL. On the stored per-head features of \ladNRuns{} runs, averaging the five metrics over heads after computing them per head scores \ladMetricAvg{}, keeping them apart \ladPerHead{}, a gain of \ladGainPerHead{} that is positive on \ladGainPerHeadPos{} of \ladNRuns{} runs; permuting the head axis independently for every item, which keeps every value and destroys only which head produced it, costs \ladGainVsShuffle{}, and padding the averaged features with noise to the per-head width changes them by \ladNoiseEffect{}. The signal is in which head does what.

The reason is elementary and we state it once. Any functional of the head-averaged attention graph is constant on head configurations that share a mean, so no averaged read can represent the contrast between heads (\cref{prop:annihilation}). For the algebraic connectivity in particular, the averaged graph can only overstate it, since the second Laplacian eigenvalue is a pointwise minimum of linear functionals and hence concave, with equality exactly when every head shares a Fiedler vector (\cref{prop:jensen}); on the layers we measured the averaged graph reports connectivity well above the mean over heads throughout. LapEigvals already keeps heads apart, and its authors note that the causal Laplacian is triangular so its eigenvalues are its diagonal \citep{binkowski2025lapeigvals}; those features are therefore an in-degree sequence, invariant to any rewiring that preserves degrees, while a symmetrised spectrum is a global functional of the routing, which is why two such different readouts perform alike once heads are kept apart (\cref{app:proofs}). Concurrent anonymous work proves that neither can represent which token attends to which, and that a readout anchored on a distinguished row can; on tool calls the anchored readout ties the per-head profile (\GapAnchPerHeadLlamaOneBBfcl{} and \GapAnchPerHeadGemmaBfcl{}), so the limit it identifies is not what separates attention from the residual stream here.
